\begin{table*}[p]
\scriptsize
\centering
\caption{\textbf{Setting-specific covariate sets for the Section~2.3 adjusted models.} Each model was fitted separately within the six task settings. Task ecology is therefore fixed within a setting and is not interpreted as a within-setting coefficient. The table lists the estimable covariates used to adjust the scaffolded-support association; pooled context models use task ecology separately as shown in Supplementary Tables~\ref{tab:constructive_context_logit} and \ref{tab:constructive_rate_context_logit}.}
\label{tab:setting_model_covariates}
\setlength{\tabcolsep}{3pt}
\renewcommand{\arraystretch}{1.12}
\begin{tabularx}{\textwidth}{>{\raggedright\arraybackslash}p{0.13\textwidth}>{\raggedright\arraybackslash}p{0.28\textwidth}>{\raggedright\arraybackslash}p{0.22\textwidth}>{\raggedright\arraybackslash}X}
\toprule
Settings & Primary Poisson count model & Logistic any-constructive model & Sensitivities and notes \\
\midrule
WC coding; LMSYS coding; SC coding; WC writing; LMSYS writing; SC writing
& Scaffolded-support presence; user framing; total turns; emotional-expression ratio; observable error marker; persistence-after-failure marker; high-persistence marker.
& Scaffolded-support presence; user framing; total turns.
& The offset-rate sensitivity used the same Poisson mean model but represented user turns with an exposure offset. User-framing stratified Poisson models omitted user framing because it was fixed within stratum. Error and persistence markers are most interpretable in coding-oriented contexts but are retained as observed conversation-level context when present. \\
\bottomrule
\end{tabularx}
\end{table*}
